\chapter{Where to go from here}

\section{What you have now}

Part I took you from an empty \lstinline|main()| to a first person shooter. Part
II walked the rest of the library module by module. Between the two, the code in
this repository calls about 240 of raylib's 619 functions --- and the rest is
mostly depth inside modules you have already met, not new territory.

More to the point: you can now read \file{raylib.h} and understand what you are
looking at. That was always the real goal, because that header is the
documentation and nothing else is.

\section{The examples are the next teacher}

The official examples live in
\file{build/external/raylib-master/examples}, 225 programs sorted by module. With
this manual behind you they are readable rather than intimidating.

Worth opening first, in this order:

\begin{center}
\begin{tabular}{@{}ll@{}}
\toprule
\file{core/core\_3d\_camera\_first\_person.c} & the official version of chapter 12 \\
\file{models/models\_first\_person\_maze.c} & the official version of chapter 14 \\
\file{shapes/shapes\_collision\_area.c} & collision resolution, done cleanly \\
\file{textures/textures\_sprite\_anim.c} & sprite sheets in practice \\
\file{audio/audio\_sound\_multi.c} & overlapping sounds without clipping \\
\file{shaders/shaders\_basic\_lighting.c} & where the flat look finally goes away \\
\file{others/rlgl\_standalone.c} & the layer under raylib, when you need it \\
\bottomrule
\end{tabular}
\end{center}

\section{Three things this manual did not cover}

\begin{enumerate}[leftmargin=*,itemsep=5pt]
  \item \textbf{raygui.} An immediate-mode GUI library in a single header, and it
        is \emph{already in your examples folder}
        (\file{examples/*/raygui.h}). Buttons, sliders, checkboxes, dropdowns and
        text boxes, in the same style as the rest of raylib:
\begin{lstlisting}
#define RAYGUI_IMPLEMENTATION
#include "raygui.h"

if (GuiButton((Rectangle){ 20, 20, 120, 30 }, "Start")) StartGame();
GuiSlider((Rectangle){ 20, 60, 200, 20 }, "vol", NULL, &volume, 0.0f, 1.0f);
\end{lstlisting}
        For a pause menu or an options screen it turns a day of work into an hour.

  \item \textbf{rlgl.} The layer raylib itself is built on: a thin abstraction
        over OpenGL that lets you push vertices directly. You want it the day you
        need a custom blend mode, instanced rendering or geometry raylib has no
        function for. It is \file{src/rlgl.h}, and it is standalone.

  \item \textbf{Shipping.} Building for the web with Emscripten, packaging for
        Windows from Linux, bundling assets. The project wiki on raylib.com is
        the reference, and the examples folder ships a working
        \file{Makefile.Web} to start from.
\end{enumerate}

\section{The tools on raylib.com}

Ramon Santamaria, raylib's author, also writes small tools that pair with it.
They are not part of the library and you download them separately:

\begin{center}
\begin{tabular}{@{}ll@{}}
\toprule
\textbf{rFXGen} & generates retro sound effects with sliders \\
\textbf{rTexPacker} & packs sprites into an atlas \\
\textbf{rTexViewer} & views and converts textures \\
\textbf{rGuiStyler} & themes for raygui \\
\textbf{rGuiLayout} & lays out raygui screens and exports the C code \\
\textbf{rIconPacker} & builds \file{.ico} files \\
\bottomrule
\end{tabular}
\end{center}

rFXGen in particular is worth ten minutes: chapter 15 showed you how those sounds
are synthesised, and the tool lets you find one by ear instead of by arithmetic.

\section{Habits worth keeping}

\begin{itemize}[leftmargin=*,itemsep=4pt]
  \item Update and draw stay separate. Update changes variables, draw paints them.
  \item Every speed is per second and gets multiplied by \lstinline|dt|.
  \item Every \lstinline|Load...| has an \lstinline|Unload...|, outside the loop,
        on both sides.
  \item Detect collisions with one call; resolve them one axis at a time.
  \item Anything inside a \lstinline|Begin|/\lstinline|End| pair is in that
        pair's space. The HUD goes outside.
  \item When something is invisible, check the draw order and the camera before
        suspecting anything clever.
  \item When a collision misbehaves, draw the box.
  \item Game state in a struct, plus a reset function.
  \item When stuck, open \file{raylib.h} and search. The answer is one grep away.
\end{itemize}

\section{A last word}

The gap between this manual and a game you would show someone is not more
knowledge --- it is repetitions. You have the loop, the cameras, the collisions,
the maths, the sound and the shaders. What is left is choosing something small,
finishing it, and discovering the twenty tiny problems that only show up when you
try.

Make the mini Doom yours. Add a door that opens. Add a key that opens it. Add a
sound when it does. Add a second level. Each of those is an afternoon, and each
one teaches more than a chapter.
